\documentclass[12pt]{article}

\input{../../_assets/preambulo.tex}


\begin{document}

    % 1. Foto de fondo
    % 2. Título
    % 3. Encabezado Izquierdo
    % 4. Color de fondo
    % 5. Coord x del titulo
    % 6. Coord y del titulo
    % 7. Fecha

    
    \input{../../_assets/portada}
    \portadaExamen{ffccA4.jpg}{Análisis Funcional\\Examen XVI}{Análisis Funcional. Examen XVI}{MidnightBlue}{-8}{28}{2025/26}{David Sánchez Muñoz}

    \begin{description}
        \item[Asignatura] Análisis Funcional.
        \item[Curso Académico] 2025/26.
        \item[Grado] Grado en Informática y Matemáticas.
        \item[Grupo] Único.
        \item[Profesor] David Arcoya Álvarez.
        \item[Descripción] Examen final extraordinario.
        \item[Fecha] 12 de febrero del 2026.
        % \item[Duración] ---.
    
    \end{description}
    \newpage


    % ------------------------------------
    
    \begin{ejercicio}[3 puntos]
        Sea $\{y_n\}_{n \in \bb{N}} \subset \bb{R}$ tal que
        \[ \sum_{n \geq 1} x_n y_n \text{ converge para todo } \{x_n\}_{n \in \bb{N}} \in \ell^2. \]
        Prueba que $\{y_n\}_{n \in \bb{N}} \in \ell^2$.
    \end{ejercicio}

    \begin{ejercicio}[3 puntos]
        Sean $(X, \|\cdot\|_X)$ y $(Y, \|\cdot\|_Y)$ dos espacios de Banach y $T : X \longrightarrow Y$ lineal y acotada para la que existe $c > 0$ tal que
        \[ \|Tx\|_Y \geq c\|x\|_X, \quad \forall x \in X. \]
        Prueba que $T$ es compacto si y solo si $\dim X < \infty$.
    \end{ejercicio}

    \begin{ejercicio}[4 puntos]
        Sean $E = \ell^p$ y $F = \ell^q$ con $p, q \in \left]1, \infty\right[$. Suponiendo que $a : \bb{R} \longrightarrow \bb{R}$ es una función continua verificando
        \[ |a(t)| \leq C|t|^{\nicefrac{p}{q}}, \quad \forall t \in \bb{R}, \]
        se define la aplicación $A : \ell^p \longrightarrow \ell^q$ mediante
        \[ Ax = \{a(x(k))\}_{k \in \bb{N}}, \quad \text{para cada } x = \{x(k)\}_{k \in \bb{N}} \in \ell^p. \]
        \begin{enumerate}[label=\textbf{\alph*)}]
            \item \textbf{[1.5 puntos]} Prueba que $A$ está bien definida y es continua de $\ell^p$ (con la topología de la norma) en $\ell^q$ (con la topología de la norma).
            \item \textbf{[1.5 puntos]} Demuestra que si $\{x^n\}_{n \in \bb{N}}$ es una sucesión en $\ell^p$ que converge ($\sigma(E, E^*)$) débilmente a $x \in \ell^p$, entonces $\{A(x^n)\}_{n \in \bb{N}}$ converge ($\sigma(F, F^*)$) débilmente a $Ax \in \ell^q$.
            \item \textbf{[1 punto]} Deduce que la restricción de $A$ a la bola unidad cerrada $\ol{B}_E$ es continua de $(\ol{B}_E, \sigma(E, E^*))$ en $(F, \sigma(F, F^*))$.
        \end{enumerate}
    \end{ejercicio}

    \newpage
    \setcounter{ejercicio}{0} % Reiniciar contador de ejercicios
    \noindent

    % El alumno redactará aquí la solución.
\end{document}